\documentclass[11pt]{article}

% Change "review" to "final" to generate the final (sometimes called camera-ready) version.
% Change to "preprint" to generate a non-anonymous version with page numbers.
\usepackage[preprint]{acl}

% Standard package includes
\usepackage{times}
\usepackage{latexsym}

% For proper rendering and hyphenation of words containing Latin characters (including in bib files)
\usepackage[T1]{fontenc}
% For Vietnamese characters
\usepackage[T5]{fontenc}
% See https://www.latex-project.org/help/documentation/encguide.pdf for other character sets

% This assumes your files are encoded as UTF8
\usepackage[utf8]{inputenc}

% This is not strictly necessary, and may be commented out,
% but it will improve the layout of the manuscript,
% and will typically save some space.
\usepackage{microtype}

% This is also not strictly necessary, and may be commented out.
% However, it will improve the aesthetics of text in
% the typewriter font.
\usepackage{inconsolata}

%Including images in your LaTeX document requires adding
%additional package(s)
\usepackage{graphicx}
\usepackage{booktabs}
\usepackage{multirow}
\usepackage{amsmath}
\usepackage{amsfonts}
\usepackage{bm}

\title{Multi-Aspect Heterogeneous Graph Neural Networks for \\ Vietnamese News Recommendation}

\author{T. Nhat, \ L. M. Nhut, \ N. H. Phat \\
  Faculty of Computer Science \\
  University of Information Technology \\
  \texttt{\{2152xxxx, 2152xxxx, 2152xxxx\}@gm.uit.edu.vn} \\\And
  Msc. Huynh Van Tin \\
  Faculty of Information Science and Technology \\
  University of Information Technology \\
  \texttt{tinhv@uit.edu.vn} \\}

\begin{document}
\maketitle

\begin{abstract}
  News recommendation in Vietnamese presents unique challenges due to linguistic complexity, sparse user interactions, and the "cold-start" problem. In this work, we propose a comprehensive breakdown of the recommendation pipeline, leveraging Multi-Aspect Heterogeneous Graph Neural Networks (MA-HGN) and Contrastive Learning (CL). We introduce a novel heterogeneous graph construction incorporating users, articles, categories, and authors. We benchmark state-of-the-art Collaborative Filtering (CF) models, including LightGCL, SimGCL, and XSimGCL, alongside our custom implementations of MA-HCL and HGN. Our extensive experiments on a real-world VnExpress dataset utilizing 15,693 users and 3,645 articles demonstrate that while Content-Based methods (PhoBERT) still lead in global retrieval (Recall@10 of 0.0535), our HGN model (G2) achieves superior stability in Leave-One-Out ranking (Recall@10 of 0.4364), efficiently capturing social homophily signals.
\end{abstract}

\section{Introduction}

In the era of information overload, personalized news recommendation systems (NRS) are essential for helping users discover relevant content. While significant progress has been made in English-language NRS, Vietnamese NRS remains under-explored, primarily due to the lack of large-scale datasets and the complexity of the language.

Traditional approaches often rely on collaborative filtering (CF), which suffers from data sparsity, or content-based (CB) methods, which may lack serendipity. Graph Neural Networks (GNNs) have emerged as a powerful tool to bridge this gap, capable of modeling complex interactions between users and items. However, most existing GNN-based NRS model the data as a bipartite user-item graph, ignoring rich auxiliary information such as article categories, authors, and temporal dynamics.

In this paper, we address these limitations by modeling the news ecosystem as a \textbf{Heterogeneous Information Network (HIN)}. Our key contributions are:
\begin{enumerate}
  \item \textbf{Heterogeneous Graph Construction:} We construct extensive graph variants (G1-G4), integrating Author and Category nodes, and introducing weighted edges based on user engagement (comments, replies) and article freshness.
  \item \textbf{Advanced Modeling:} We implement and adapt \textbf{SimMAHGN} (Simple Multi-Aspect Heterogeneous Graph Network) and \textbf{MA-HCL} (Multi-Aspect Heterogeneous Contrastive Learning) to the Vietnamese context.
  \item \textbf{Comprehensive Benchmarking:} We provide a rigorous evaluation of various embeddings (VN-SBERT, PhoBERT, BGE-M3) and CF models (SimGCL, XSimGCL, LightGCL).
  \item \textbf{Deployment:} We develop an interactive Streamlit application demonstrating real-time hybrid recommendation with A/B testing capabilities.
\end{enumerate}

\section{Related Work}

\subsection{Graph Neural Networks for Recommendation}
Early graph-based recommendation methods like NGCF \citet{wang2019neural} explicitly modeled high-order connectivities in the user-item graph. However, the heavy computational cost of non-linear activations led to the proposal of LightGCL \citet{he2020lightgcn}, which simplified the GCN architecture by removing linear transformations and non-linearities. Recently, SimGCL \citet{yu2022graph} integrated contrastive learning (CL) into LightGCL, using random noise perturbations to enforce representation consistency, significantly improving robustness against sparse data.

\subsection{Heterogeneous Graph Recommendation}
Heterogeneous GNNs (HGNNs) extend GNNs to graphs with multiple node and edge types. Methods like HAN and HGT utilize meta-paths to aggregate neighbors. In news recommendation, \citet{liu2021heterogeneous} proposed modeling news as a heterogeneous graph including entities and categories, which inspired our MA-HCL approach. Unlike previous works that focus solely on English datasets (MIND), we adapt these architectures for Vietnamese news by integrating language-specific embeddings.

\subsection{Pretrained Language Models in SRS}
PLMs like BERT \citet{devlin2018bert} have revolutionized content understanding. In the Vietnamese context, PhoBERT and VN-SBERT \citet{reimers2019sentence} provide strong baselines. Our work systematically evaluates these embeddings within a hybrid NRS framework.

\section{Methodology}

\subsection{Data Collection and Preprocessing}
We crawled data from VnExpress, a major Vietnamese news portal, resulting in 3,645 articles and 40,045 user interactions. The data includes:
\begin{itemize}
  \item \textbf{Users (15,693):} Anonymized user IDs, usernames, and join dates.
  \item \textbf{Articles (3,645):} Title, short description, content, publication date, and author.
  \item \textbf{Interactions (40,045):} User comments and replies, serving as a proxy for engagement. Calculated sparsity is approx. 99.93\%.
\end{itemize}

\subsection{Heterogeneous Graph Construction}
We define a heterogeneous graph $G = (V, E)$ where $V = \mathcal{U} \cup \mathcal{I} \cup \mathcal{C} \cup \mathcal{A}$ consists of user, article, category, and author nodes.
We experimented with three primary graph tiers to isolate the impact of different signals:
\begin{enumerate}
  \item \textbf{G1 - Bipartite Baseline:} Standard User-Article graph ($E_{ui}$) with comment interactions.
  \item \textbf{G2 - Social Expansion:} Adds User-User edges ($E_{uu}$) based on reply interactions, capturing social homophily. Total edges increase by approx. 61k.
  \item \textbf{G3 - Semantic Expansion:} Adds Category nodes ($E_{uc}, E_{ci}$), hypothesizing that categories act as semantic hubs to reduce path length.
\end{enumerate}

\subsection{Model Architecture}

\subsubsection{SimGCL \& XSimGCL}
We adopted SimGCL \citet{yu2022graph} as a strong baseline foundation. SimGCL simplifies Graph Convolutional Networks (GCN) by removing non-linearities and introducing contrastive learning (CL) noise to regularize embedding representations.

Let $E^{(k)}$ be the embedding matrix at layer $k$. The propagation rule is:
\begin{equation}
  E^{(k)} = \frac{1}{2} (D^{-\frac{1}{2}} A D^{-\frac{1}{2}}) E^{(k-1)}
\end{equation}
SimGCL augments the representation learning with a contrastive auxiliary task:
\begin{equation}
  \mathcal{L}_{CL} = \sum_{i \in \mathcal{U} \cup \mathcal{I}} - \log \frac{\exp(sim(e_i^{'}, e_i^{''}) / \tau)}{\sum_{j \neq i} \exp(sim(e_i^{'}, e_j^{''}) / \tau)}
\end{equation}
where $e_i^{'}$ and $e_i^{''}$ are perturbed embeddings with added random noise.

\subsubsection{MA-HCL (Multi-Aspect Heterogenous CL)}
Adapted from \citet{liu2021heterogeneous}, MA-HCL handles heterogeneity using meta-path based aggregation. For each meta-path $\Phi$ (e.g., User-Article-Author-Article), we compute a distinct embedding $h_u^{\Phi}$ for user $u$.
We employ a \textbf{hierarchical attention mechanism} to fuse embeddings from different meta-paths:
\begin{equation}
  h_u = \sum_{\Phi} \alpha_{\Phi} h_u^{\Phi}
\end{equation}
where $\alpha_{\Phi}$ is the learned importance of meta-path $\Phi$. To improve robustness, we apply contrastive learning on the meta-path views, encouraging consistency between the global representation and aspect-specific representations.

\subsubsection{Hybrid Inference Strategy}
To address the cold-start problem where graph connectivity is sparse or non-existent (new users/items), we implement a hybrid strategy.
Let $s_{CF}(u, i)$ be the score from the CF model (SimGCL/MA-HCL).
Let $s_{CB}(u, i)$ be the cosine similarity between user history and candidate item embedding (VN-SBERT).
The final score is:
\begin{equation}
  s_{final} = \beta \cdot \sigma(s_{CF}) + (1-\beta) \cdot s_{CB}
\end{equation}
where $\beta$ is a dynamic weight parameter regulated by user activity level.

\section{Experiments}

\subsection{Experimental Setup}
\textbf{Dataset Split:} We use two evaluation protocols:
\begin{enumerate}
  \item \textbf{Full Ranking (Global):} 80/20 random split. Ranks all ~3.6k articles to simulate cold-start discovery.
  \item \textbf{Leave-One-Out (Loo):} Ranks 1 relevant item against 99 random negatives. Focuses on user-specific preference stability.
\end{enumerate}

\textbf{Baselines:} LightGCN \cite{he2020lightgcn}, SimGCL \cite{yu2022graph}, LightGCL \cite{cai2023lightgcl}, XSimGCL \cite{yu2023xsimgcl}, and our HGN (HGRec-based) and HCL (Contrastive).

\subsection{Main Results}
Table \ref{tab:results_full} and \ref{tab:results_loo} summarize performance.

\begin{table}[h]
  \centering
  \small
  \caption{Benchmark: Full Ranking (Recovering items from the entire corpus).}
  \label{tab:results_full}
  \setlength{\tabcolsep}{3pt}
  \begin{tabular}{lccc}
    \hline
    \textbf{Model} & \textbf{Graph} & \textbf{Recall@10} & \textbf{NDCG@10} \\
    \hline
    PhoBERT        & CB             & \textbf{0.0535}    & \textbf{0.0315}  \\
    TF-IDF         & CB             & 0.0508             & 0.0301           \\
    \hline
    LightGCL       & G2             & 0.0436             & 0.0261           \\
    XSimGCL        & G2             & 0.0432             & 0.0260           \\
    HGN (Ours)     & G2             & 0.0369             & 0.0201           \\
    HCL (Ours)     & G2             & 0.0320             & 0.0192           \\
    SimGCL         & G2             & 0.0228             & 0.0131           \\
    \hline
  \end{tabular}
\end{table}

\begin{table}[h]
  \centering
  \small
  \caption{Benchmark: Leave-One-Out (Precision vs Noise).}
  \label{tab:results_loo}
  \setlength{\tabcolsep}{3pt}
  \begin{tabular}{lccc}
    \hline
    \textbf{Model}      & \textbf{Graph} & \textbf{Recall@10} & \textbf{NDCG@10} \\
    \hline
    \textbf{HGN (Ours)} & \textbf{G2}    & \textbf{0.4364}    & \textbf{0.2636}  \\
    LightGCL            & G1             & 0.3081             & 0.2021           \\
    XSimGCL             & G2             & 0.3036             & 0.2002           \\
    HCL (Ours)          & G3             & 0.2260             & 0.1492           \\
    SimGCL              & G3             & 0.2181             & 0.1333           \\
    \hline
  \end{tabular}
\end{table}

In **Full Ranking**, Content-Based methods (PhoBERT) currently lead, indicating that for pure exploration, semantic content matching is reliable. However, in **Leave-One-Out**, our **HGN** model significantly outperforms all baselines (R@10 0.4364 vs 0.3081 for LightGCL), proving that when the search space is constrained, social signals (G2) provide highly accurate specific recommendations.

\subsection{Ablation Studies}

\subsubsection{Impact of Graph Structure}
We compared G1 (Bipartite), G2 (Social), and G3 (Category).
\begin{itemize}
  \item \textbf{G1 (Base):} Density 0.06\%, Recall@10 0.0262.
  \item \textbf{G2 (Social):} Density 0.064\%, Recall@10 \textbf{0.0369} (+40\% gain). Social edges provide high-quality homophily signals.
  \item \textbf{G3 (Category):} Density 0.096\%, Recall@10 0.0299. Interestingly, adding category hubs *decreased* performance compared to G2. This suggests that category edges may be too generic, diluting the specific user interests ("Information Overload" at a structural level).
\end{itemize}

\section{Conclusion}
We presented a comprehensive study of Graph Neural Networks for Vietnamese news recommendation. Our work underscores the value of heterogeneous graphs, specifically the inclusion of social reply networks (G2). While global retrieval remains challenging due to extreme sparsity (99.93\%), our HGN model demonstrates state-of-the-art stability in personalized ranking tasks.

\section*{Limitations}
Our dataset is relatively small (15k users). The "Full Ranking" performance (Recall $\sim$ 3-5\%) highlights the difficulty of purely structural learning in ultra-sparse environments without hybrid content injection. GNN training on denser graphs (G2/G3) also faces scalability challenges, with MA-HGN taking 2-3x longer than XSimGCL.

\bibliography{custom}

\end{document}
